%% Use the "normalphoto" option if you want a normal photo instead of cropped to a circle
% \documentclass[10pt,a4paper,normalphoto]{altacv}

\documentclass[9pt,letterpaper,ragged2e,withhyper]{../asset/altacv}
%% AltaCV uses the fontawesome5 and packages.
%% See http://texdoc.net/pkg/fontawesome5 for full list of symbols.

% Change the page layout if you need to
\geometry{left=1cm,right=1cm,top=0.6cm,bottom=0.7cm,columnsep=0.8cm}

% The paracol package lets you typeset columns of text in parallel
\usepackage{paracol}
\usepackage{hyphenat}

% Change the font if you want to, depending on whether
% you're using pdflatex or xelatex/lualatex
\iftutex
  % If using xelatex or lualatex:
  \setmainfont{Roboto Slab}
  \setsansfont{Lato}
  \renewcommand{\familydefault}{\sfdefault}
\else
  % If using pdflatex:
  \usepackage[rm]{roboto}
  \usepackage[defaultsans]{lato}
  % \usepackage{sourcesanspro}
  \renewcommand{\familydefault}{\sfdefault}
\fi

% Change the colours if you want to
\definecolor{SlateGrey}{HTML}{2E2E2E}
\definecolor{LightGrey}{HTML}{666666}
\definecolor{DarkPastelRed}{HTML}{450808}
\definecolor{PastelRed}{HTML}{8F0D0D}
\definecolor{GoldenEarth}{HTML}{E7D192}
\colorlet{name}{black}
\colorlet{tagline}{PastelRed}
\colorlet{heading}{DarkPastelRed}
\colorlet{headingrule}{GoldenEarth}
\colorlet{subheading}{PastelRed}
\colorlet{accent}{PastelRed}
\colorlet{emphasis}{SlateGrey}
\colorlet{body}{LightGrey}



% Change some fonts, if necessary
\renewcommand{\namefont}{\Huge\rmfamily\bfseries}
\renewcommand{\personalinfofont}{\footnotesize}
\renewcommand{\cvsectionfont}{\Large\rmfamily\bfseries}
\renewcommand{\cvsubsectionfont}{\large\bfseries}


% Change the bullets for itemize and rating marker
% for \cvskill if you want to
\renewcommand{\itemmarker}{{\small\textbullet}}
\renewcommand{\ratingmarker}{\faCircle}
\newcommand{\compactdivider}{%
  \par\noindent\textcolor{body!30}{\hdashrule{\linewidth}{0.4pt}{0.5ex}}\par
}

%% Use (and optionally edit if necessary) this .tex if you
%% want to use an author-year reference style like APA(6)
%% for your publication list
% \input{pubs-authoryear}

%% Use (and optionally edit if necessary) this .tex if you
%% want an originally numerical reference style like IEEE
%% for your publication list
% \input{pubs-num}

%% main.bib contains your publications
% \addbibresource{main.bib}

\hyphenpenalty=1000
\exhyphenpenalty=1000

\begin{document}
\name{Sam (Shengkai) Xu}
\tagline{Ph.D. Researcher | Machine Learning, Computer Vision \& Engineering Systems}
%% You can add multiple photos on the left or right
\photoR{2cm}{../asset/shengkai}
% \photoL{2.5cm}{Yacht_High,Suitcase_High}

\personalinfo{%
  % Not all of these are required!
  \email{shengkai.x.sam@gmail.com}
  \email{sxu7@uncc.edu}
  \phone{734-731-6398}
  \location{Charlotte, NC}
  \linkedin{shengkai-xu-sam}
  \github{samxu29}
  %% You can add your own arbitrary detail with
  %% \printinfo{symbol}{detail}[optional hyperlink prefix]
  % \printinfo{\faPaw}{Hey ho!}[https://example.com/]
  %% Or you can declare your own field with
  %% \NewInfoFiled{fieldname}{symbol}[optional hyperlink prefix] and use it:
  % \NewInfoField{gitlab}{\faGitlab}[https://gitlab.com/]
  % \gitlab{your_id}
  %%
  %% For services and platforms like Mastodon where there isn't a
  %% straightforward relation between the user ID/nickname and the hyperlink,
  %% you can use \printinfo directly e.g.
  % \printinfo{\faMastodon}{@username@instace}[https://instance.url/@username]
  %% But if you absolutely want to create new dedicated info fields for
  %% such platforms, then use \NewInfoField* with a star:
  % \NewInfoField*{mastodon}{\faMastodon}
  %% then you can use \mastodon, with TWO arguments where the 2nd argument is
  %% the full hyperlink.
  % \mastodon{@username@instance}{https://instance.url/@username}
}

\makecvheader
%% Depending on your tastes, you may want to make fonts of itemize environments slightly smaller
% \AtBeginEnvironment{itemize}{\small}

%% Set the left/right column width ratio to 6:4.
\columnratio{0.66}

% Start a 2-column paracol. Both the left and right columns will automatically
% break across pages if things get too long.
\begin{paracol}{2}

\cvsection{Experience}

\cvevent{Research Assistant}{UNC Charlotte, Dr.\ Hongfei Xue}{2024 -- Present}{}
\begin{itemize}
\item Architect, train, and evaluate deep-learning and computer-vision models for multimodal sensing with mmWave radar, RGB/depth cameras, LiDAR, and acoustic sensors.
\item Build ground-truth and sensor-data pipelines for pose and mesh reconstruction; analyze accuracy, latency, and cross-modal consistency and communicate results to cross-disciplinary collaborators.
\end{itemize}

\divider

\cvevent{AI Research Assistant}{UNC Charlotte TeCSAR Lab}{2022 -- 2023}{}
\begin{itemize}
\item Developed a video-processing and anomaly-action-detection pipeline for civilian security and public-safety applications.
\item Increased processing throughput by 50\%, improving video smoothness and system response time.
\end{itemize}

\divider

\cvevent{Embedded Firmware Engineer}{Oxit LLC}{2021 -- 2023}{}
\begin{itemize}
\item Developed modular Embedded C libraries across chipsets and built low-power RF and IoT systems.
\item Designed AWS and Google Cloud automation pipelines and tools that improved development efficiency and compatibility.
\end{itemize}

\cvsection{Selected Projects}

\textbf{Neural Musculoskeletal Pose Estimation}\\
\textit{Human Mesh Reconstruction, Inverse Kinematics | 2026--Present}\\
\begin{itemize}
\item Real-time mesh and pose pipeline being translated toward deployment with collaborators in sports medicine.
\end{itemize}

\compactdivider
\textbf{SocialBit On-Device Classifier}\\
\textit{Wearable ML, Audio Classification | 2026--Present}\\
\begin{itemize}
\item Privacy-focused smartwatch models for interaction type, conversational depth, and tone under edge-compute constraints.
\end{itemize}

\cvsection{Selected Publications}

\textbf{Real-Time Neural Musculoskeletal Pose Estimation}\\
\textit{IEEE/ACM CHASE 2026}
\begin{itemize}
\item Replaced iterative inverse kinematics with a neural surrogate in a real-time monocular reconstruction pipeline.
\end{itemize}

\divider

\textbf{Towards Robust mmWave-Based Human Activity Recognition Using Large Simulated Datasets for Model Pretraining}\\
\textit{IEEE BigData 2024}
\begin{itemize}
\item Used simulated human-mesh data and heatmap perturbations to improve pretraining with limited real-world data.
\end{itemize}

\divider

\textbf{Real-Time Object Detection Algorithm Performance on Edge Computing Devices}\\
\textit{IEEE ISNCC 2024}
\begin{itemize}
\item Benchmarked YOLOv5 latency and power on Jetson Nano, Raspberry Pi 4, and Coral TPU.
\end{itemize}

\switchcolumn
\fontsize{10pt}{12pt}\selectfont

\cvsection{Profile}

\begin{quote}
Ph.D. researcher developing deep-learning models and experimental pipelines for physical sensing, edge deployment, and engineering prediction. Embedded-systems and CAD experience; Formula SAE contributor motivated to apply AI and scientific computing to motorsports.
\end{quote}

\vfill
\cvsection{Education}

\cvevent{Ph.D.\ Computer Science}{UNC Charlotte}{Expected 2029}{}
\small Research focus: deep-learning systems for multimodal sensing.

\divider

\cvevent{M.Sc.\ Computer Engineering}{University of North Carolina \\ at Charlotte}{Aug 2022 -- Aug 2024}{}
\small Research focus: machine learning on edge and embedded systems.

\divider

\cvevent{B.Sc.\ Computer Engineering}{University of North Carolina \\ at Charlotte}{Jan 2018 -- Dec 2022}{}
\small Machine Learning concentration; Mathematics minor

\vfill
\cvsection{Technical Skills}

\textbf{Programming \& ML}\\
Python, C++, Embedded C, SQL, PyTorch, TensorFlow, OpenCV, scikit-learn

\compactdivider
\textbf{Engineering \& Systems}\\
SOLIDWORKS, KiCad, sensor integration, experimental data acquisition, embedded systems, Linux, Git/GitHub, AWS, Google Cloud

\vfill
\cvsection{Motorsports Involvement}

\textbf{Formula SAE Electrical Team | 2018--2019}\\
Contributed to electrical power distribution and vehicle wiring for UNC Charlotte Formula SAE.

\compactdivider
\textbf{Club-Level Automobile Time Attack | 2020--2024}\\
Participated in time-attack events and prepared personal track cars through mechanical maintenance, repairs, inspections, and track-day setup.

\compactdivider
\textbf{Motorcycle Track Days | 2024--Present}\\
Participate in track days and prepare personal motorcycles through routine service, mechanical repairs, and pre-event safety inspections.

\end{paracol}

% \newpage


% % \cvsection{Solutions Intelligence Developer} % Cover-Letter
% \makecvheader


% \bigskip
% \divider

% \begin{figure}[h]
%   \centering
%   \begin{minipage}{0.3\textwidth}
%   \centering
%   \includegraphics[width=\linewidth]{qr_IEEE.png}
%   \end{minipage}\hfill
%   \begin{minipage}{0.3\textwidth}
%   \centering
%   \includegraphics[width=\linewidth]{qr_github.png}
%   \end{minipage}\hfill
%   \begin{minipage}{0.3\textwidth}
%   \centering
%   \includegraphics[width=\linewidth]{qr_linkedin.png}
%   \end{minipage}
%   \end{figure}

% \today

% \bigskip

% Dear Hiring Manager,

% \bigskip
% I am writing to express my strong interest in the Software Engineering Summer Scholar position at Deloitte Consulting. I am currently pursuing my Master's degree in Computer Engineering at the University of North Carolina at Charlotte, with an expected graduation date of May 2024. As evidenced by my attached resume, I believe I possess the skills, experience, and motivation to thrive in this role.

% \bigskip
% As a firmware engineer intern at Oxit LLC, I gained hands-on experience developing embedded systems and IoT applications using technologies like LoRa and low-power RF communication. In this role, I designed module libraries to streamline embedded development across different chipsets. I also served as an AI research assistant at the UNC Charlotte TeCSAR lab, where I worked on an anomaly detection pipeline for public safety video analytics. These experiences allowed me to hone my software development and machine learning skills.

% \bigskip
% In addition to my work experience, I have undertaken several programming projects on my own initiative. For instance, I developed a Telegram-based automated 3D printing server for remote printing. I also created a real-time mask detection system using YOLO object detection models. Through projects like these, I have become proficient in languages like Python, C++, and PyTorch/TensorFlow machine learning frameworks.

% \bigskip
% I believe that my academic coursework has also prepared me well for this Software Engineering role. I have taken classes covering software engineering principles, data structures and algorithms, operating systems, and computer architecture. I am confident in my ability to apply this knowledge to develop robust and scalable cloud-native applications using agile methodologies.

% \bigskip
% I have also applied for your Business Technology Solutions Summer Scholar roles. I believe that my passion for learning new technologies paired with my engineering background enables me to thrive in a consulting environment, understand business needs, and deliver technological solutions. I am excited by the opportunity to gain well-rounded experience through a rotational program like Deloitte's.

% \bigskip
% In summary, I am eager to leverage my technical skills and continual drive to learn to deliver impactful solutions on diverse client engagements. I would welcome the opportunity to discuss how I can contribute to Deloitte Consulting as a Software Engineering or Business Technology Solutions Summer Scholar. Thank you for your consideration, and I look forward to hearing from you.

% \bigskip
% Sincerely, \\
% Shenkai Xu

\end{document}
